\documentclass[a4paper,11pt]{article}

\usepackage[english,greek]{babel}
\usepackage[utf8]{inputenc}
\usepackage[T1]{fontenc}
\usepackage{amsmath}
\usepackage{amssymb}
\usepackage{graphicx}
\usepackage{booktabs}
\usepackage{float}
\usepackage{subcaption}
\usepackage{hyperref}
\usepackage{listings}
\usepackage{xcolor}
\usepackage[left=2.5cm,right=2.5cm,top=2.5cm,bottom=2.5cm]{geometry}

% Ρυθμίσεις για την εμφάνιση κώδικα
\definecolor{codegreen}{rgb}{0,0.6,0}
\definecolor{codegray}{rgb}{0.5,0.5,0.5}
\definecolor{codepurple}{rgb}{0.58,0,0.82}
\definecolor{backcolour}{rgb}{0.95,0.95,0.92}

\lstdefinestyle{mystyle}{
    backgroundcolor=\color{backcolour},   
    commentstyle=\color{codegreen},
    keywordstyle=\color{magenta},
    numberstyle=\tiny\color{codegray},
    stringstyle=\color{codepurple},
    basicstyle=\ttfamily\footnotesize,
    breakatwhitespace=false,         
    breaklines=true,                 
    captionpos=b,                    
    keepspaces=true,                 
    numbers=left,                    
    numbersep=5pt,                  
    showspaces=false,                
    showstringspaces=false,
    showtabs=false,                  
    tabsize=2
}

\lstset{style=mystyle}

\newcommand{\lt}[1]{\foreignlanguage{english}{#1}}

\title{Άσκηση: \lt{KPCA + LDA} \\ Υλοποίηση και Αποτελέσματα}
\author{Όνοματεπώνυμο: Γεώργιος Κεσογλίδης \\ ΑΕΜ: 220}
\date{\today}
\begin{document}

\maketitle

\begin{abstract}
Στην παρούσα εργασία παρουσιάζεται η υλοποίηση και η πειραματική αξιολόγηση ενός συνδυασμού μη-γραμμικής μείωσης διαστάσεων (\lt{Kernel PCA}) και (\lt{Linear Discriminant Analysis -- LDA}) για \lt{multi-class image classification}. Χρησιμοποιήθηκαν υποσύνολα του \lt{CIFAR-10}. Η υλοποίηση και οι πειραματικές ρουτίνες έγιναν σε \lt{Python} (\lt{NumPy} / \lt{PyTorch}) μέσα σε \lt{Jupyter Notebook} και τα αποτελέσματα αποθηκεύτηκαν ως εικόνες και πίνακες για την τελική αναφορά.
\end{abstract}

\tableofcontents
\newpage

\section{Εισαγωγή}
Στόχος της εργασίας είναι να αξιολογηθεί κατά πόσο η μη-γραμμική μείωση διαστάσεων μέσω του \lt{KPCA} βελτιώνει τη διαχωριστική ικανότητα όταν συνδυαστεί με την \lt{LDA} σε προβλήματα εικόνων (π.χ. \lt{CIFAR-10}). Το κείμενο καλύπτει τη μαθηματική θεμελίωση, την πρακτική υλοποίηση, την πειραματική διαδικασία και την ανάλυση των αποτελεσμάτων.

\section{Μεθοδολογία και Υλοποίηση}
\subsection{\lt{Kernel PCA} (\lt{KPCA})}
Χρησιμοποιήθηκε ο \lt{RBF} πυρήνας:
\begin{equation}
k(\mathbf{x},\mathbf{y}) = \exp(-\gamma \|\mathbf{x}-\mathbf{y}\|^2).
\end{equation}
Στην υλοποίηση του \lt{notebook} υπάρχει η κλάση \lt{KPCA} με τα εξής κρίσιμα σημεία:
\begin{itemize}
    \item Υπολογισμός και κεντράρισμα του πίνακα πυρήνα $K$ με τον \lt{RBF} πυρήνα, αποθηκεύοντας τα μέσα για το κεντράρισμα νέων δειγμάτων.
    \item Ιδιοτιμητική αποσύνθεση του κεντραρισμένου $K_c$ (\lt{scipy.linalg.eigh}), με φιλτράρισμα πολύ μικρών/αρνητικών ιδιοτιμών (παράμετρος \texttt{tol}).
    \item Κανονικοποίηση ιδιοδιανυσμάτων (διαιρώντας με $\sqrt{\lambda}$) και αποθήκευση των επιλεγμένων ιδιοτιμών/ιδιοδιανυσμάτων.
\end{itemize}

\subsection{\lt{Linear Discriminant Analysis} (\lt{LDA})}
Η \lt{LDA} υπολογίζει τους ενδο- ($S_W$) και δια- ($S_B$) διασπορικούς πίνακες και λύνει το γενικευμένο ιδιοτιμητικό πρόβλημα
\begin{equation}
S_B \mathbf{w} = \lambda S_W \mathbf{w}.
\end{equation}
Στην παρούσα εργασία χρησιμοποιήθηκε μια απλή υλοποίηση (ορίζεται ως \lt{LDAClassifier} στο \lt{notebook}) με τα παρακάτω βασικά βήματα:
\begin{enumerate}
    \item Υπολογισμός μέσων ανά κλάση $\mu_c$, συνολικού μέσου $\mu$, και των πινάκων $S_W$ και $S_B$.
    \item Λύση του γενικευμένου ιδιοτιμητικού προβλήματος (\lt{scipy.linalg.eigh}) και επιλογή των κορυφαίων συνιστωσών \lt{($\text{n\_components}=\min(\text{n\_features},\text{n\_classes}-1)$).}
    \item Κατασκευή πίνακα προβολής $W$ και προ-υπολογισμός προβαλλόμενων μέσων κλάσεων $\mu_c^\top W$.
    \item Ταξινόμηση νέων δειγμάτων με βάση το πλησιέστερο προβαλλόμενο μέσο \lt{(nearest-class-mean).}
\end{enumerate}

\newpage

\section{Πειραματική Διαδικασία}
\subsection{Δεδομένα}
Χρησιμοποιήθηκε το σύνολο \lt{CIFAR-10}. Για αποφυγή μεγάλου υπολογιστικού κόστους χρησιμοποιήθηκαν υποσύνολα (π.χ. 2000 δείγματα εκπαίδευσης, 1000 δοκιμής) για την εκπαίδευση του \lt{KPCA}. Έγινε κανονικοποίηση ανά κανάλι (mean/std) και \lt{flatten} για τον υπολογισμό του πυρήνα.

\subsection{Γραφήματα από \lt{KPCA + LDA}}
\begin{figure}[H]
    \centering
    \begin{subfigure}{0.48\textwidth}
        \centering
            \includegraphics[width=\textwidth]{results/kpca_eigen_spectrum.png}
            \caption{\lt{Eigenvalue spectrum of KPCA}}
    \end{subfigure}\hfill
    \begin{subfigure}{0.48\textwidth}
        \centering
        \includegraphics[width=\textwidth]{results/kpca_2d_train.png}
            \caption{\lt{KPCA: first two components (train)}}
    \end{subfigure}
    \vspace{3mm}
    \begin{subfigure}{0.48\textwidth}
        \centering
        \includegraphics[width=\textwidth]{results/lda_2d_train.png}
            \caption{\lt{LDA projection (train)}}
    \end{subfigure}\hfill
    \begin{subfigure}{0.48\textwidth}
        \centering
        \includegraphics[width=\textwidth]{results/cifar10_kpca_lda_scatter.png}
            \caption{\lt{KPCA+LDA test projection (scatter)}}
    \end{subfigure}
    \caption{Κύρια γραφήματα από τα πειράματα}
    \label{fig:all_plots}
\end{figure}

\subsection{Σύγκριση με διαφορετικούς κατηγοριοποιητές}
Παρακάτω παρουσιάζονται συνοπτικά τα αριθμητικά αποτελέσματα από την αξιολόγηση διαφορετικών κατηγοριοποιητών (\lt{accuracy}, \lt{balanced accuracy}, \lt{macro precision/recall/F1}, χρόνοι εκπαίδευσης/πρόβλεψης).

% Σύντομες περιγραφές μοντέλων (μία γραμμή έκαστο)
\noindent\textbf{Σύντομες περιγραφές των αξιολογηθέντων μοντέλων:}
\begin{itemize}
    \item \lt{KPCA + LDA (custom): RBF-KPCA (n\_components=100,} $\gamma=1e-5$) με \lt{custom LDAClassifier} με \lt{nearest-class-mean.}
    \item \lt{kNN (KPCA): k-NN (k=3)} πάνω στις \lt{KPCA} αναπαραστάσεις.
    \item \lt{Logistic (KPCA)}: Πολυωνυμική/πολυκατηγορική \lt{Logistic Regression} πάνω στις \lt{KPCA} συνιστώσες.
    \item \lt{RandomForest (KPCA)}: Σύνολο δέντρων απόφασης \lt{(100 estimators)} πάνω στο KPCA.
    \item \lt{PCA + LDA (custom)}: Γραμμική \lt{PCA (n\_components=100)} και στη συνέχεια η ίδια custom \lt{LDAClassifier}.
    \item \lt{kNN (PCA): k-NN (k=3)} πάνω στο\lt{PCA}.
    \item \lt{SVM (PCA)}: \lt{SVM} με \lt{RBF kernel} πάνω στο \lt{PCA}.
\end{itemize}


\begin{table}[H]
\centering
\caption{Σύνοψη αποτελεσμάτων στο \lt{CIFAR-10} (υποσύνολα)}
\label{tab:results}
\begin{tabular}{lccccccc}
	\toprule
	\textbf{Μοντέλο} & \textbf{\lt{Acc}} & \textbf{\lt{BalAcc}} & \textbf{\lt{Prec$_{mac}$}} & \textbf{\lt{Rec$_{mac}$}} & \textbf{\lt{F1$_{mac}$}} & \textbf{\lt{Fit (s)}} & \textbf{\lt{Pred (s)}} \\
\midrule
\lt{KPCA + LDA (custom)} & 0.415 & 0.417 & 0.417 & 0.417 & 0.416 & 0.013 & 0.012 \\
\lt{RandomForest} (\lt{KPCA}) & 0.324 & 0.325 & 0.332 & 0.325 & 0.320 & 0.433 & 0.049 \\
\lt{Logistic} (\lt{KPCA}) & 0.304 & 0.308 & 0.343 & 0.308 & 0.274 & 0.048 & 0.001 \\
\lt{kNN} (\lt{KPCA}) & 0.247 & 0.249 & 0.350 & 0.249 & 0.238 & 0.001 & 0.011 \\
\lt{PCA + LDA (custom)} & 0.107 & 0.107 & 0.107 & 0.107 & 0.106 & 0.009 & 0.011 \\
\lt{SVM} (\lt{PCA}) & 0.109 & 0.109 & 0.105 & 0.109 & 0.094 & 0.427 & 0.332 \\
\lt{kNN} (\lt{PCA}) & 0.108 & 0.113 & 0.100 & 0.113 & 0.089 & 0.001 & 0.020 \\
\bottomrule
\end{tabular}
\end{table}

\begin{figure}[H]
    \centering
    \begin{subfigure}{0.48\textwidth}
        \centering
        \includegraphics[width=\textwidth]{results/cifar10_evaluation_metrics.png}
        \caption{\lt{Metrics comparison}}
    \end{subfigure}\hfill
    \begin{subfigure}{0.48\textwidth}
        \centering
        \includegraphics[width=\textwidth]{results/cifar10_best_confusion_norm.png}
        \caption{\lt{Normalized confusion matrix (best model)}}
    \end{subfigure}
    \caption{Σύγκριση μοντέλων και \lt{confusion matrix}}
\end{figure}


\subsection{Σύγκριση με \lt{scikit-learn}}
Επιπλέον, έγινε άμεση σύγκριση της χειροποίητης υλοποίησης με τις υλοποιήσεις του \lt{scikit-learn} για \lt{KPCA} + \lt{LDA}.

\begin{figure}[H]
    \centering
    \begin{subfigure}{0.48\textwidth}
        \centering
        \includegraphics[width=\textwidth]{results/cifar10_kpca_lda_confusion.png}
        \caption{\lt{KPCA+LDA confusion matrix (test)}}
    \end{subfigure} \hfill 
    \begin{subfigure}{0.48\textwidth}
        \centering
        \includegraphics[width=\textwidth]{results/kpca_lda_sklearn_confusion.png}
        \caption{\lt{Confusion matrix - sklearn KPCA+LDA}}
    \end{subfigure}
    \caption{\lt{KPCA+LDA: test set visualizations}}
    \label{fig:kpca_test_vis}
\end{figure}


\begin{figure}[H]
    \centering
    \begin{subfigure}{1.0\textwidth}
        \centering
        \includegraphics[width=\textwidth]{results/kpca_lda_comparison_scatter.png}
        \caption{\lt{Custom vs sklearn KPCA+LDA projections}}
    \end{subfigure}
    \caption{Σύγκριση της \lt{custom} υλοποίησης με τις υλοποιήσεις του \lt{scikit-learn}}
    \label{fig:sklearn_compare}
\end{figure}

\section{Συμπεράσματα}
Από τα παραχθέντα γραφήματα (\lt{eigen‑spectrum}, \lt{2‑D} προβολές, \lt{confusion matrices}, ανά‑κλάση) εξάγονται τα εξής συμπεράσματα:
\begin{itemize}
    \item Το \lt{KPCA+LDA} υπερέχει του \lt{PCA+LDA} στο υποσύνολο $(\approx 41.5\%$ ακρίβεια) με βελτίωση \lt{macro‑F1} και \lt{balanced accuracy.}
    \item Το \lt{eigen‑spectrum} δείχνει ότι οι πρώτες συνιστώσες συγκεντρώνουν το κύριο σήμα. Η επιλογή \lt{n\_components}=100 χρειάζεται έλεγχο σε μεγαλύτερα δείγματα.
    \item Συγχέσεις εντοπίζονται σε οπτικά παρόμοιες κλάσεις (π.χ. \lt{cat/dog}). Για βελτίωση απαιτούνται πλουσιότερα χαρακτηριστικά ή περισσότερα δείγματα.
\end{itemize}

\end{document}